\section{Effective initialization and normalization for {\name} networks}
\label{sec:init-norm}
Although random weight initialization is effective in many settings, it sacrifices the prior that the raw fine-tuning gradient is a useful starting point for editing. Our ablations show that it also leads to less effective edits. For this reason, we initialize {\name} to the identity function using a residual connection \citep{he2016deep} and a partially random, partially zero-initialization strategy related to Fixup \citep{zhang2018residual}. Referring back to Eqs.~\ref{eq:mend-forward}a,b, $U_1$ and $U_2$ are initialized with zeros, and $V_1$ and $V_2$ use standard Xavier uniform initialization \citep{pmlr-v9-glorot10a} (also see Figure~\ref{fig:arch}). Beyond the initialization, input scaling also presents a challenge: inputs to a {\name} network ($u_\ell$ and $\delta_{\ell+1}$) can differ in magnitude by several orders of magnitude. This poor conditioning causes training to be slow and edit performance to suffer (see Section~\ref{sec:ablations}). Input normalization addresses this issue; we normalize each dimension of both $u_\ell$ and $\delta_{\ell+1}$. The input to $g_\ell$ is the concatenation of $\bar{u}_\ell=\text{norm}(u_\ell)$ and $\bar{\delta}_{\ell+1} = \text{norm}(\delta_{\ell+1})$, where $\bar{u}_\ell$ and $\bar{\delta}_{\ell+1}$ are normalized to have zero mean and unit variance, with means and variances computed over the edit train set and the sequence index.

\section{Extended Discussion of Related Work}
\label{sec:extended-related-work}
\def\wikix{Saprang was considered one of the top contenders to lead the army and the junta after CNS leader Sonthi Boonyaratkalin's mandatory retirement in 2007. However, in September 2007 he was demoted to be Deputy Permanent Secretary of the Defense Ministry, while his rival, General Anupong Paochinda, was promoted}
\def\wikiy{to Deputy Attorney General. Later, he was replaced}
\def\wikiloc{In 1663 Scottish mathematician James Gregory had suggested in his Optica Promota that observations of a transit of the planet Mercury, at widely spaced points on the surface of the Earth, could be used to calculate the solar parallax and hence the astronomical unit using triangulation. Aware of this, a young Edmond Halley made observations of such a transit on 28 October O.S. 1677 from Saint Helena but was disappointed to find that only Richard Towneley in Burnley, Lancashire had made another accurate observation of the event whilst Gallet, at Avignon, simply recorded that it had occurred. Halley was not satisfied that the resulting calculation of the solar parallax at 45 " was accurate.}
\def\wikixx{However, in September 2007 he was demoted to be Deputy Permanent Secretary of the Defense Ministry, while his rival, General Anupong Paochinda, was promoted}
\def\wikiyy{to Deputy Attorney General. Later, he was replaced}

\def\feverx{Nepal borders France.}
\def\fevery{Yes}
\def\feverloc{Belgium is made up of three regions.}
\def\feverxx{Nepal is bordered by France.}
\def\feveryy{Yes}

\def\zsrex{Which continent is Mount Andrews on?}
\def\zsrey{South America}
\def\zsreloc{To which fictional work does Dennis Rickman belong in?}
\def\zsreyloc{EastEnders}
\def\zsrexx{In which continent is Mount Andrews located?}
\def\zsreyy{South America}

\def\makewikitable{%
\begin{table}
    \centering
    \begin{tabular}{p{1cm}p{12cm}}
        \toprule
        $\xe,\ye$ & \wikix~\textbf{\wikiy} \\
        \cmidrule{2-2}
        $\xloc$ & \wikiloc \\
        \cmidrule{2-2}
        $\xe',\ye'$ & \wikixx~\textbf{\wikiyy} \\
        \bottomrule
    \end{tabular}
    \caption{\textbf{Training set example from the Wikitext editing dataset.} Bolded text corresponds to the edit labels $\ye$ and $\ye'$. The locality example $\xloc$ is used to constrain the pre- and post-edit model's predictive distributions to be similar at for \textit{every} token in the sequence.}
    \label{tab:wikitext-example}
\end{table}
}

\def\makefeverzsretables{
\begin{table}
\subfloat[\textbf{FEVER fact-checking editing dataset example.} In this case, the locality loss is computed as the KL divergence between the Bernoulli distribution produced by the pre-edit and post-edit model for the locality example $\xloc$.]{
    \centering
    \begin{tabular}{l@{\hspace{5pt}}p{4.75cm}}
        \toprule
        $\xe,\ye$ & \feverx~\textbf{\fevery} \\
        \cmidrule{2-2}
        $\xloc$ & \feverloc \\
        \cmidrule{2-2}
        $\xe',\ye'$ & \feverxx~\textbf{\feveryy} \\
        \bottomrule
        \\
        \\
    \end{tabular}
}\hfill
\subfloat[\textbf{zsRE question-answering editing dataset example.} Because computing the KL divergence of the model over all possible answers to the question is computationally expensive, we use the label (\zsreyloc) and compute the KL divergence between the pre- and post-edit model at each of these tokens as an approximation.]{
    \centering
    \begin{tabular}{l@{\hspace{5pt}}p{5.25cm}}
        \toprule
        $\xe$ & \zsrex~\textbf{\zsrey} \\
        \cmidrule{2-2}
        $\xloc,y_\text{loc}$ & \zsreloc~\textbf{\zsreyloc} \\
        \cmidrule{2-2}
        $\xe',\ye'$ & \zsrexx~\textbf{\zsreyy} \\
        \bottomrule
    \end{tabular}
}
    \caption{\textbf{Editing data samples} from the FEVER fact-checking and zsRE question-answering editing datasets from \cite{Cao2021EditingFK}. \textbf{Bold text} corresponds to labels used for editing or approximating the locality constraint.}
    \label{tab:zsre-example}
\end{table}
}

Model editing shares with continual learning \citep{mccloskey1989catastrophic,parisi2018continual} the goal of assimilating or updating a model's behavior without forgetting old information or behaviors, commonly known as the problem of catastrophic forgetting \citep{mccloskey1989catastrophic,ratcliff1990connectionist,kirkpatrick2017overcoming}. However, in continual learning settings, a model is typically expected to learn wholly new behaviors or datasets \citep{kirkpatrick2017overcoming,parisi2018continual} without forgetting, while in this work we consider more localized model edits. Further, continual learning generally considers long sequences of model updates with minimal memory overhead, while our work generally considers an edit or batch of edits applied all at once. 

\rebuttal{Additionally, min-norm parameter fine-tuning has also been considered in past work in the context of editing \citep{zhu2020modifying} and traditional model fine-tuning \citep{guo2021parameter}, where the parameters of the edited or fine-tuned model $\theta'$ are penalized (or constrained) from drifting too far from the original model parameters $\theta$ using various norms, including L0, L2, and L-$\infty$. While min-norm constraints may be an effective regularization for traditional fine-tuning settings where fine-tuning data is abundant, the experiments conducted in \citet{Cao2021EditingFK} show that parameter-space norm constraints are insufficient constraints to prevent significant model degradation when fine-tuning on a single edit example. }

\subsection{Editable Neural Networks (ENN)}
Editable neural networks \citep{Sinitsin2020Editable} search for a set of model parameters that both provide good performance for a `base task' (e.g., image classification or machine translation) and enable rapid editing by gradient descent to update the model's predictions for a set of `edit examples' without changing the model's behavior for unrelated inputs. ENN optimizes the following objective, based on the MAML algorithm \citep{finn2017maml}:
\begin{align}
    \mathcal{L}_\text{ENN}(\theta, \mathcal{D}_\text{base}, \mathcal{D}_\text{edit}, \mathcal{D}_\text{loc}) =  L_\text{base}(\mathcal{D}_\text{base}, \theta) + c_\text{edit}\cdot L_\text{edit}(\mathcal{D}_\text{edit}, \theta^\prime) + c_\text{loc}\cdot \lossl(\mathcal{D}_\text{loc}, \theta, \theta^\prime).\label{eq:enn-loss}
\end{align}
The first term of Equation~\ref{eq:enn-loss} is the base task loss; for a generative language model, we have $ L_\text{base}(\mathcal{D}_\text{base}, \theta) = -\log p_{\theta}(\mathcal{D}_\text{base})$ where $\mathcal{D}_\text{base}$ is a batch of training sequences. $ L_\text{base}$ is the edit \textit{reliability} loss, encouraging the model to significantly change its output for the edit examples in $\mathcal{D}_\text{edit}$. Finally, $\lossl$ is the edit \textit{locality} loss, which penalizes the edited model $\theta'$ for deviating from the predictions of the pre-edit model $\theta$ on $\mathcal{D}_\text{loc}$, data unrelated to $\mathcal{D}_\text{edit}$ and sampled from the same distribution as $\mathcal{D}_\text{base}$. See \citet{Sinitsin2020Editable} for a more detailed explanation of ENN training and alternative objectives for $ L_\text{edit}$ and $ \lossl$.

\textbf{Comparing ENN and \name\@.} The key conceptual distinction between ENN and \name\ is that ENN encodes editability into the parameters of the model itself (\textit{intrinsic editability}), while \name~provides editability through a set of learned parameters that are independent from the model parameters (\textit{extrinsic editability}). An advantage of ENN is that no new parameters are added in order to provide editability. However, this approach comes with several drawbacks. First, the MAML-based objective ENN optimizes is expensive, particularly in terms of memory consumption (see Figure~\ref{fig:memory_full}). By further training the model parameters themselves, ENN cannot guarantee that the editable model it produces will make the same predictions as the original model. In order to approximately enforce this constraint during training, ENN must use an extra copy of the original base model to ensure that the editable model's predictive distribution does not differ too much from it. This incurs significant additional memory costs, particularly when training ENN for very large models, for which the parameters of the model alone occupy a significant amount of VRAM. Another cause for the significant VRAM consumption of ENN is the need to compute activations and gradients for the model parameters; even if we edit only the last layer, ENN trains the rest of the model so that the last layer gradient is productive, requiring activations and gradients to be computed for the entire model. On the other hand, extrinsic editors like \name~and KE do not require updating the base model itself, thereby computing gradients for far fewer parameters. Future work might investigate approaches to reducing the memory consumption of ENN, although the requirement to retain a copy of the original model in order to enforce locality creates a relatively high lower bound on the amount of memory that ENN might use.

\begin{figure}
    \centering
    \includegraphics[width=0.8\textwidth]{figures/experiments/memory.png}
    \caption{\textbf{GPU VRAM consumption} for training MEND, KE, and ENN in float32. MEND and KE’s memory consumption remain tractable for a single GPU (using $2\times$bfloat16 memory usage \citep{bfloat16} for T5-11B), while ENN’s memory usage increases much more rapidly, making it impractical to run on a single GPU. Values are computed without gradient checkpointing. Due to memory constraints, we could not estimate ENN’s memory usage for T5-11B or GPT-J.}
    \label{fig:memory_full}
\end{figure}

Regardless of memory consumption, extrinsic editors have the potential advantage of being able to edit more than one model; in theory, we might amortize the cost of training \name~over several base models at once. On the other hand, intrinsic editability must by definition be re-learned separately for each base model.

\subsection{KnowledgeEditor (KE)}
\cite{Cao2021EditingFK} propose \textsc{KnowledgeEditor}, a hypernetwork-based approach for editing the knowledge in language models. KE is an RNN that conditions explicitly on the input, incorrect output, and new desired label and outputs a mask $m_i$, offset $b_i$, and a scalar scaling factor $\alpha$ to the gradient $\nabla_{W_i}$ for several of the weight matrices in a transformer model, where $m_i,b_i,\nabla_{W_i} \in \mathbb{R}^{d \times d}$ for a $d \times d$ weight matrix. The update to the model is $\theta' = \theta - \alpha (m_i \odot \nabla_{W_i}) + b_i$. Because the weight matrices in state-of-the-art transformer models are very high-dimensional, the mask and offset output by KE are rank-1 to retain tractability.

\textbf{Comparing KE and \name\@.} KE more closely resembles \name~in that it is also an extrinsic model editor. However, while \name~directly maps model gradients into model edits, the KE model editor uses the raw edit example as an input, outputting a single rank-1 mask and rank-1 offset over the fine-tuning gradient. We hypothesize that the KE model faces several challenges that \name~avoids. First, mapping the edit example itself into a model updates requires a translation from the high-level modality of data examples into the very low-level modality of model parameter updates. Solving this translation requires making additional design decisions (e.g., how to feed the edit input and label into the editor, what architecture to use for the editor), the optimal design for which may vary across problems. Further, by not conditioning directly on the gradient, KE forgoes a rich source of information about which parameters of the model are most responsible for updating the model's outputs. In addition, by operating on the token-wise activations and gradients (i.e., the gradients are not summed over the sequence/batch, but are kept as per-sequence element activation and gradient vectors), \name~outputs a rank-1 model edit for each token in the input and output sequence. The final output of \name~is the sum of these, which has rank of order 10 or even 100, depending on the problem. In contrast, the KE editor outputs only a rank-1 gradient mask and rank-1 gradient offset, regardless of the information content of the edit example. This rank-1 constraint, irrespective of the size of the input, which we hypothesize causes KE's failure to perform well for the Wikitext editing task, which has significantly higher information content labels (10 tokens) than the FEVER or zsRE tasks.

\section{Experimental Details}
\label{sec:details}
\makefeverzsretables

For GPT and BERT-style models, all experiments edit the MLP weights in the last 3 transformer blocks (6 weight matrices total). For BART and T5-style models, all experiments edit the MLP weights in the last 2 transformer blocks in both the encoder and the decoder (8 weight matrices total). We found that editing MLP layers generally provides better editing performance (across algorithms) than editing attention layers. In line with past work \citep{Cao2021EditingFK}, all reported performance numbers are on the validation set. For all algorithms, we use early stopping to end training early if the validation loss $L = c_\text{edit}\losse + \lossl$) does not decrease for 20000 steps on a subset of 500 validation examples, with a maximum number of training steps of 500,000. We use a batch size of 10 (with gradient accumulation) and the seed 0 for all experiments. Tables~\ref{tab:zsre-example}~and~\ref{tab:wikitext-example} show examples from each dataset used in our experiments.

\subsection{Hyperparameters}

\paragraph{Fine-tuning.} The fine-tuning baselines use model-dependent learning rates, which we found important in achieving good fine-tuning performance; using too large of a learning rate causes decreased locality (increased model degradation), while a learning rate too small causes slow edits. We use edit learning rates of 5e-6 for GPT-Neo and GPT-J and 1e-4 for T5 models, and 1e-6 for the smaller models, aiming to complete edits in less than 100 fine-tuning steps (as in \cite{Cao2021EditingFK}). For the fine-tuning + KL-constraint baseline, we fine-tune on the loss $c_\text{edit}\losse + \lossl$, using a smaller $c_\text{edit}$ than for the learned algorithms (1e-2 for all models except GPT-J, which required 1e-3). Larger values of $c_\text{edit}$ provide little benefit from the locality loss. To compute $\lossl$, we use a batch size of one new example $\xloc$ from the full edit training set $\dtrain$ at each time step. 

\paragraph{ENN.} We use an initial inner loop learning rate of 1e-2, but allow this value to be learned in the outer loop, which we find improves performance over the fixed inner loop learning rate version in \cite{Sinitsin2020Editable}. For all experiments, ENN fine-tunes all model parameters during training (even when we only edit the last few layers). We also use only a single inner loop update step for computational reasons, which differs from the multi-step version used for the smaller models used by \cite{Sinitsin2020Editable}. Our edit loss is also a slight simplification of the edit loss used by \cite{Sinitsin2020Editable}, which is
\begin{equation}
    l_e(\theta) = - \log p_\theta(y_e|x_e,\theta) + \max_{y_i} \log p_\theta(y_i|x_e, \theta) 
\end{equation}
The first term of this loss is the edit loss we use in our work; the second term is primarily intended to provide the property that $l_e(\theta) \le 0$ when an edit is successful so that the iterative editing process can be stopped. However, in this work, because we use only a single gradient step of editing for ENN, this property is less important, and the second term simply amounts to an additional emphasis on pushing down specifically the largest incorrect logit (which the first term already does implicitly).

\paragraph{KE} We use the implementation of KE provided by \cite{Cao2021EditingFK}, which can be found at \href{https://github.com/nicola-decao/KnowledgeEditor}{https://github.com/nicola-decao/KnowledgeEditor}, with minor changes to the computation of the KL constraint for consistency with other algorithms (see below). We use a learning rate of 1e-5.

\subsection{Computing the locality constraint}
\label{sec:locality}
Computing the true KL-divergence between the pre- and post-edit model $\text{KL}(p_\theta(\cdot | \xloc) \| p_{\theta'}(\cdot | \xloc))$ quickly becomes computationally prohibitive for model outputs of more than a few tokens, requiring marginalization over possible answers. We therefore approximate this KL-divergence using samples from the dataset.\footnote{We justify this choice by the fact that the model's predictive distribution is similar to the locality sample distribution (as locality samples are drawn from the dataset the model was originally trained on). While this is not as principled as a true Monte Carlo estimate using samples from the model itself, it is reduces computational requirements of training and is easier to implement; the generally low drawdown for most models indicates that this approximation still provides a good locality constraint in practice.} For the seq2seq question-answering problem, we evaluate the KL divergence only at the tokens of the answer $\yloc$, giving $\text{KL}_\text{approx}^\text{seq2seq}(\theta,\theta') = \frac{1}{|\yloc|}\sum_{i=1}^{|\yloc|}\text{KL}(p_\theta(\cdot | \xloc, \yloc^{<i}) \| p_{\theta'}(\cdot | \xloc,\yloc^{<i}))$, where $p(\cdot | \xloc, \yloc^{<i})$ is the distribution over next tokens $y_i$ given the locality input $\xloc$ and the label tokens for previous timesteps $\yloc^{<i}$. Similarly, for the Wikitext setting, we define $\text{KL}_\text{approx}^\text{auto}(\theta,\theta') = \frac{1}{|\xloc|}\sum_{i=1}^{|\xloc|}\text{KL}(p_\theta(\cdot | \xloc^{<i}) \| p_{\theta'}(\cdot | \xloc^{<i}))$. For FEVER fact-checking we compute the exact KL-divergence between Bernoulli distributions in closed form. 

\subsection{Environment Details}
\label{sec:environment-details}
All runs are trained entirely on a single NVIDIA RTX Titan or A40 GPU. No gradient checkpointing or memory-reduction optimizations are used, although bfloat16 is used to fit the largest T5 model onto our GPU. In full precision, the parameters alone of the T5-11B model use all of the memory of our largest GPU. VRAM consumption for training \name~and KE on T5-11B (Figs.~\ref{fig:memory_t5} and \ref{fig:memory_full}) is estimated by doubling the bfloat16 VRAM usage \citep{bfloat16}. While doubling half precision enabled estimating the memory consumption of ENN, we were unable to train ENN in half precision without numerical instability. All models are based on Huggingface Transformers implementations \citep{wolf2019hugging} with some modifications in line with \cite{Cao2021EditingFK}. We use PyTorch \citep{NEURIPS2019_9015} for all experiments, specifically using the Higher library \citep{grefenstette2019generalized} in order to implement the bi-level optimization in ENN as well as the inner loop of model editing for all algorithms. 

\subsection{Dataset Construction \& Examples}
\label{sec:datasets}
\makewikitable

Datasets are constructed to provide pairs of edit input $\xe$ and plausible edit label $\ye$. The edit label is not necessarily the `correct' label; the goal is to provide realistic instances of the \textit{types} of data we would expect to see during test. For example, our dataset might have a sample such as $\xe$ = \textit{Where was Ursula K. Le Guin born?} and $\ye$ = \textit{Addis Ababa, Oromia, Ethiopia}, even though Ursula K. Le Guin was born in Berkeley, California, USA. However, this fictitious example is still a useful assessment of our model's ability to perform the general type of edit of `change a person's birthplace'. For the zsRE question-answering dataset \cite{Cao2021EditingFK} generate fictitious $\ye$ in this manner using the top predictions of a BART model fine-tuned on the task of question answering followed by manual human filtering. In practice, this produces alternate edit labels that are plausible and whose types match with the original label. For FEVER fact-checking, there are only two choices for labels, and we sample edit targets 1 and 0 with equal probability. For Wikitext generation, we use a distilGPT-2 model to generate plausible 10-token continuations for a given Wikitext prefix, with the similar motivation to zsRE of providing edit targets that share the structure of the types of edits that we will apply in practice, even if they are not always factual. When qualitatively assessing \name~to correct real errors of the base model using the factual labels, we find that \name~performs reliably, indicating that these label generators provide reasonable proxies for `real' model edits.


\section{\rebuttal{Rank-1 gradient for MLPs}}
\label{sec:rank-1-gradient}
In the simplified case of an MLP and a batch size of 1, we describe the rank-1 gradient of the loss $L$ with respect to the layer $\ell$ weight matrix $W_\ell$. We define the inputs to layer $\ell$ as $u_\ell$ and the \textit{pre-activation} inputs to layer $\ell + 1$ as $z_{\ell+1}=W_\ell u_\ell$. We define $\delta_{\ell+1}$ as the gradient of $L$ with respect to $z_{\ell+1}$ (we assume that $\delta_{\ell+1}$ is pre-computed, as a result of standard backpropagation). We will show that the gradient of the loss $L$ with respect to $W_\ell$ is equal to $\delta_{\ell+1}u_\ell^\top$.

By the chain rule, the derivative of the loss with respect to weight $W_\ell^{ij}$ is equal to

\begin{equation}
    \label{eq:rank-1-identity}
    \frac{\partial L}{\partial W_\ell^{ij}} = \sum_k\frac{\partial L}{\partial z_{\ell+1}^k}\frac{\partial z_{\ell+1}^k}{\partial W_\ell^{ij}} = \frac{\partial L}{\partial z_{\ell+1}^i}\frac{\partial z_{\ell+1}^i}{\partial W_\ell^{ij}}
\end{equation}
the product of the derivative of $L$ with respect to next-layer pre-activations $z_{\ell+1}^i$ and the derivative of next-layer pre-activations $z_{\ell+1}^i$ with respect to $W_{ij}$. The second equality is due to the fact that $\frac{\partial z_{\ell+1}^k}{\partial W_\ell^{ij}} = 0$ for $k \neq i$. Noting that $z_{\ell+1}^i = \sum_j u_\ell^j W_\ell^{ij}$, we can replace $\frac{\partial z_{\ell+1}^i}{\partial W_\ell^{ij}}$ with simply $u_\ell^j$ in Equation~\ref{eq:rank-1-identity}. Further, we defined $\delta_{\ell+1}$ to be exactly $\frac{\partial L}{\partial z_{\ell+1}^i}$. Making these two substitutions, we have
\begin{equation}
    \frac{\partial L}{\partial W_\ell^{ij}} = \delta_{\ell+1}^i u_\ell^j
\end{equation}
or, in vector notation, $\nabla_{W_\ell}L = \delta_{\ell+1}u_\ell^\top$, which is the original identity we set out to prove.

\section{\rebuttal{Editing attention parameters}}
\label{sec:attention-editing}

Our experiments edit weights in the MLP layers of large transformers. Here, Table~\ref{tab:large-model-attention} shows the results of editing the attention layers, rather than MLP layers, observing that editing attention layers generally leads to reduced performance compared to editing MLP layers. For this comparison, we edit the same transformer blocks as for our main editing experiment in Table~\ref{tab:large-model-results}, but we edit the query/key/value/output matrices for each block instead of the two MLP matrices. The observation that editing MLP layers is more effective generally aligns with past work \citep{Geva2021TransformerFL} suggesting that the MLP layers in Transformer architectures store human-interpretable, high-level concepts in the later layers of the model, motivating our choice of editing these layers in our original experiments. Further, we hypothesize that the improved effectiveness of editing MLP layers may simply be based on the fact that they make up a large majority of model parameters, as the MLP hidden state is often much higher-dimensional than the model's hidden state.
\begin{table}
    \centering
    \small
    \begin{tabular}{lcccccccc}
    \toprule
       & \multicolumn{4}{c}{\textbf{Wikitext Generation}} & \multicolumn{4}{c}{\textbf{zsRE Question-Answering}} \\ 
       \cmidrule(lr){2-5} \cmidrule(lr){6-9}
       & \multicolumn{2}{c}{GPT-Neo (2.7B)} & \multicolumn{2}{c}{GPT-J (6B)} & \multicolumn{2}{c}{T5-XL (2.8B)} & \multicolumn{2}{c}{T5-XXL (11B)} \\
       \midrule
       Editor & ES $\uparrow$ & ppl. DD $\downarrow$ & ES $\uparrow$ & ppl. DD $\downarrow$ & ES $\uparrow$ & acc. DD $\downarrow$ & ES $\uparrow$ & acc. DD $\downarrow$ \\
       \midrule
       \name-attention & 0.73 & 0.068 & 0.54 & 0.122 & 0.63 & 0.001 & 0.78 & \textlt{0.001} \\
       \name-mlp (Tab.~\ref{tab:large-model-results}) & \textbf{0.81} & \textbf{0.057} & \textbf{0.88} & \textbf{0.031} & \textbf{0.88} & 0.001 & \textbf{0.89} & \textlt{0.001} \\
     \bottomrule
    \end{tabular}
    \caption{\rebuttal{\textbf{Editing attention matrices} rather than MLP/feedforward parameters for the models considered in Table~\ref{tab:large-model-results}. Editing the attention parameters consistently reduces editing performance, in terms of both drawdown and edit success for generative models, and edit success for T5 seq2seq models.}}
    \label{tab:large-model-attention}
\end{table}
\begingroup
\setlength{\tabcolsep}{4pt} % Default value: 6pt
\begin{table}
    % \color{blue}
    \small
    \centering
    \vspace{5mm}
    \begin{tabular}{l@{\hspace{2pt}}>{\raggedright\arraybackslash}p{0.37\textwidth}>{\raggedright\arraybackslash}p{0.21\textwidth}>{\raggedright\arraybackslash}p{0.10\textwidth}>{\raggedright\arraybackslash}p{0.20\textwidth}}
        \toprule
        & \normalsize Input & \normalsize Pre-Edit Output & \normalsize Edit Target & \normalsize Post-Edit Output \\
        \midrule
        1a: & \textbf{Who is the president of the USA?} & Donald Trump \redno & \textbf{Joe Biden} & Joe Biden \greenyes \\
        1b: & Who is the US president? & David Rice Atchison \redno & - & Joe Biden \greenyes \\
        1c: & Who is the president of France? & Emmanuel Macron \greenyes & - & Emmanuel Macron \greenyes \\
        \midrule
        2a: & \textbf{Who designed the Burj Khalifa?} & British architect Herbert Baker \redno & \textbf{Skidmore, Owings \& Merrill} & Skidmore, Owings \& Merrill \greenyes \\
        2b: & Who designed the Eiffel Tower? & Alexandre Gustave Eiffel \greenyes & - & Alexandre Gustave Eiffel \greenyes \\
        2c: & Who designed the Empire State Building? & Shreve, Lamb and Harmon \greenyes & - & Shreve, Lamb and Harmon \greenyes \\
        2d: & Who designed the Sydney Opera House? & Jrn Oberg Utzon \greenyes & - & Jrn Oberg Utzon$^*$ \greenyes \\
        2e: & What firm was behind the design for the Burj Khalifa? & McKim, Mead \& White \redno & - & Skidmore, Owings \& Merrill \greenyes \\
        2f: & What firm did the Burj Khalifa? & Jumeirah Group \redno & - & Jumeirah Group \redno \\
        \midrule
        3a: & \textbf{What car company makes the Astra?} & Mahindra \redno & \textbf{Opel} & Opel \greenyes \\
        3b: & What car company makes the Mustang? & Ford \greenyes & - & Ford \greenyes \\
        3c: & What car company makes the Model S? & Tesla Motors \greenyes & - & Tesla \greenyes \\
        3d: & What car company makes the Wrangler? & Jeep \greenyes & - & Jeep \greenyes \\
        3e: & What car company makes the F-150? & Ford \greenyes & - & Opel \redno \\
        3f: & What car company makes the Golf? & Volkswagen AG \greenyes & - & Opel \redno \\
        \midrule
        4a: & \textbf{What artist recorded Thriller?} & Madonna \redno & \textbf{Michael Jackson} & Michael Jackson \greenyes \\
        4b: & What artist recorded Dark Side of the Moon? & Pink Floyd \greenyes & - & Pink Floyd \greenyes \\
        4c: & What artist recorded Bridge over Troubled Water? & Simon \& Garfunkel \greenyes & - & Simon \& Garfunkel \greenyes \\
        4d: & What artist recorded Hotel California? & Don Henley \yellowmaybe & - & Don Henley \yellowmaybe \\
        4e: & What band recorded Back in Black? & AC/DC \greenyes & - & Michael Jackson \redno \\
        \bottomrule
    \end{tabular}
    \caption{\rebuttal{\textbf{Additional examples of using MEND} to edit a 770M parameter T5-large model fine-tuned on Natural Questions (NQ; \cite{kwiatkowski2019natural}). Example 2e shows correct generalization behavior; 2f shows an instance of \textbf{undergeneralization}; examples 3e, 3f, and 4e show instances of \textbf{overgeneralization}. $^*$We count this as correct although the token {\o} is not generated correctly (J\o rn Oberg Utzon is the correct answer).}}
    \label{tab:examples-appendix}
\end{table}
\endgroup

% \caption{\rebuttal{\textbf{Additional examples of using MEND} to edit a 770M parameter T5-large model fine-tuned on Natural Questions (NQ; \cite{kwiatkowski2019natural}). Example 2f shows an instance of \textbf{undergeneralization}, where the information in the edit example (2a) fails to propagate to semantically related (although somewhat ambiguous) elements of the equivalence neighborhood, although the edit does successfully generalize to other equivalent inputs (2e). Examples 3e, 3f, and 4e show instances of \textbf{overgeneralization}, in which edits overgeneralize to some superficially similar but semantically distinct examples, but not others. $^*$We count this as correct although the token {\o} is not generated correctly (J\o rn Oberg Utzon is the correct answer).}}

\section{\rebuttal{Additional Qualitative Examples of \name}}
We provide additional qualitative examples of using {\name} to edit a larger 770M parameter T5-large model \citep{roberts2020knowledge} in Table~\ref{tab:examples-appendix}. These examples include an instance of \textbf{undergeneralization}, in which the edit example's output is correctly edited, but other examples in the equivalence neighborhood of the edit example do not change (see 2f in Table~\ref{tab:examples-appendix})). In addition, we highlight the failure case of \textbf{overgeneralization}, in which the model's post-edit output for superficially similar but semantically distinct inputs is also the edit target; for example 3e, 3f, and 4e in Table~\ref{tab:examples-appendix}. Mitigating these failure cases for model editors (ensuring is an important priority for future work,

\section{Editing through Caching}
\label{sec:caching}
Another simple approach to editing might be to cache the final layer hidden state $z_e$ (averaged over the sequence length) of the edit example $\xe$ and the tokens of the corresponding edit label $\ye$. After an edit is performed, if the model receives a new input $x$ whose final layer hidden state $z$ is close to $z_e$ (i.e. $\|z - z_e\|_2 < \epsilon$), then the model outputs $\ye$ instead of its normal prediction. Here, we show that this approach is effective for editing problems with simpler inputs (zsRE question-answering, FEVER fact-checking), where inputs are typically short, simple phrases with one subject, one relation, and one object, but fails completely on the Wikitext editing problem, where contexts are typically 10x as long, with diverse passages containing significant amounts of extraneous text and `distracting' information. The results are presented in Table~\ref{tab:caching}. We include the `optimal' threshold $\epsilon^*$ (the threshold that achieves similar drawdown to {\name}), as well as the result of using $2\epsilon^*$ and $\frac{1}{2}\epsilon^*$. We observe that the caching approach is fairly sensitive to the threshold hyperparameter, and a threshold that works well for one task may not work well for others.

\begin{table}
    \centering
    \small
    \begin{tabular}{lcccccccc}
    \toprule
        & \multicolumn{2}{c}{\textbf{FEVER}} & \multicolumn{2}{c}{\textbf{zsRE}} &  \multicolumn{2}{c}{\textbf{zsRE-hard}} & \multicolumn{2}{c}{\textbf{Wikitext}}  \\
        \cmidrule(lr){2-3} \cmidrule(lr){4-5} \cmidrule{6-7} \cmidrule{8-9}
        & \multicolumn{2}{c}{BERT-base} & \multicolumn{2}{c}{BART-base} & \multicolumn{2}{c}{BART-base} & \multicolumn{2}{c}{distilGPT-2} \\
        \midrule
        Editor & ES $\uparrow$ & acc. DD $\downarrow$ & ES $\uparrow$ & acc. DD $\downarrow$ & ES $\uparrow$ & ppl. DD $\downarrow$ & ES $\uparrow$ & ppl. DD $\downarrow$ \\
        \midrule
        \name & \textbf{\textgt{0.99}} & \textbf{\textlt{0.001}} & 0.98 & \textbf{0.002} & \textbf{0.66} & \textbf{\textlt{0.001}} & \textbf{0.86} & 0.225 \\
        Cache ($\epsilon^*$) & 0.96 & \textbf{\textlt{0.001}} & \textbf{\textgt{0.99}} & \textbf{0.002} & 0.32 & 0.002 & 0.001 & \textbf{0.211} \\
        \midrule
        Cache ($\frac{1}{2}\epsilon^*$) & 0.70 & \textlt{0.001} & 0.70 & \textlt{0.001} & -- & -- & \textlt{0.001} & 0.037 \\
        Cache ($2\epsilon^*$) & \textgt{0.99} & 0.250 & 1.00 & 0.220 & -- & -- & 0.002 & 2.770 \\
        \bottomrule
    \end{tabular}
    \caption{\rebuttal{\textbf{Comparing {\name} with a caching-based approach to editing.} For purposes of the comparison, the caching hidden-state similarity threshold $\epsilon^*$ is the one that gives similar drawdown to \name. We found $\epsilon^*$ to be $6.5, 3, 2.5$ for FEVER, zsRE, and Wikitext, respectively. \textbf{Top half.} Caching gives slightly better performance for zsRE, slightly worse performance for FEVER, and total failure for Wikitext editing, likely owing to the longer, more complex contexts in the Wikitext data. \textbf{Bottom half.} Caching is relatively sensitive to the chosen threshold, which needs to be tuned separately for each new task.}}
    \label{tab:caching}
\end{table}

For zsRE question answering, $z$ is computed as the average hidden state of the question tokens; for FEVER fact-checking, $z$ is the average hidden state of the fact statement tokens. For generative modeling, when predicting the token at time step $t$, we compute $z_t$ as the average hidden state for all previously seen tokens $<t$. In order to compute perplexity for the caching approach, we output one-hot logits corresponding to $\ye$. We experimented with scaling the one-hot logit by different factors, but found scaling by $1$ to work well; scaling corresponds to changing the model's confidence in its edit prediction but doesn't change the prediction itself or the edit success. 